\begin{table}[htbp]\centering
\def\sym#1{\ifmmode^{#1}\else\(^{#1}\)\fi}
\caption{Top-weighted sectors: Rotemberg weights and sectoral economic size}
\label{tab:top_rotemberg_sector_size}
\small
\setlength{\tabcolsep}{3pt}
\begin{fittabular}{L{0.08\textwidth}L{0.28\textwidth}c c c c c}
\hline\hline
ISIC4 & Sector Name & Signed Rot. & Abs. Rot. & 2009 Output & Cumul. Output & China IP \\
&  & Weight & Weight & Share (\%) & Share (\%) & 2009 \\
\hline
2610 & Electronic components and boards & 0.484794 & 0.484794 & 0.098 & 0.098 & 0.6090 \\
1629 & Other wood products;articles of cork,str & 0.195351 & 0.195351 & 0.047 & 0.144 & 0.4756 \\
3290 & Other manufacturing n.e.c. & 0.157918 & 0.157918 & 0.217 & 0.362 & 0.4998 \\
2513 & Steam generators, excl. hot water boiler & 0.051172 & 0.051172 & 0.046 & 0.408 & 0.1150 \\
2431 & Casting of iron and steel & 0.049579 & 0.049579 & 0.395 & 0.802 & 0.1628 \\
\hline
\multicolumn{7}{l}{\footnotesize Notes: Sector weights decompose residualized first-stage covariance in the baseline IV sample.} \\
\multicolumn{7}{l}{\footnotesize Signed weights sum to one; absolute weights need not. Output shares are 2009 sector value added} \\
\multicolumn{7}{l}{\footnotesize relative to manufacturing. China IP 2009 is the baseline exposure share to Chinese imports.} \\
\end{fittabular}
\end{table}
